% 05_positivity_ide
\section{Positivity, and the map to interacting dark energy}
\label{sec:positivity-ide}

A recent survey of interacting dark energy models found that in seven of eight
standard couplings the energy densities become negative somewhere in the
physical domain, and it is important to quote its conclusion in full rather
than in the half that suits the present model: the negative densities arise
\emph{if energy flows from dark matter to dark energy}, which the authors read
as a strong theoretical preference for the opposite direction
\parencite{westhuizen2025_250904495,westhuizen2025_250904494,westhuizen2025_250904496}.
The model of this paper flows the disfavoured way. Its precursor is
pressureless and its product is vacuum, so if that preference is a statement
about interacting dark energy in general, this model is excluded on theoretical
grounds before any data are considered.

It is not such a statement, and this section says why. All eight couplings
surveyed have the form $Q=3H\delta\,F(\rho_{\rm dm},\rho_{\rm de})$ with $F$
depending on the recipient's density, and it is that dependence, not the
direction, which breaks positivity: a transfer whose size is set by the sector
receiving energy keeps draining a donor that has already run out. A
donor-proportional transfer in the same direction is safe, unconditionally and
with no tuning, which is what \cref{thm:quasi-positivity} establishes and what
\cref{prop:recipient-proportional} shows fails without it. The survey's
conclusion is therefore a property of the couplings it surveys and not of the
direction of flow --- a claim this paper owes the reader, having placed itself
on the wrong side of it. The section then maps the model onto the standard
parametrizations and states what it predicts that cannot be fitted away.

\subsection{Positivity is structural}\label{sec:positivity-structural}

\begin{theorem}[Quasi-positivity]\label{thm:quasi-positivity}
  Assume \cref{asm:closure} with non-negative initial data, and let $\Gam\ge0$
  be merely locally integrable along the solution: no lower bound, no upper
  bound, and no continuity are required. Then the non-negative orthant is
  forward invariant: every density in \cref{eq:cosmology-system} remains
  non-negative for as long as the solution exists.
\end{theorem}

\begin{proof}
  Not by the Nagumo subtangent criterion, which is a statement about continuous
  vector fields and would need a continuity hypothesis the theorem declines to
  make. Integrating factors suffice and need nothing beyond local
  integrability. The precursor equation is linear in $\rhoI$, so
  %
  \begin{equation*}
    \rhoI(t)=\rhoI(t_i)\exp\left[-\int_{t_i}^{t}\left(3\Hub+\Gam\right)\dd s
    \right]\ \ge 0 ,
  \end{equation*}
  %
  the integral being finite on any interval where the solution exists. With
  $\rhoI\ge0$ the sources $\bfrac\Gam\rhoI$ and $(1-\bfrac)\Gam\rhoI$ are
  non-negative, and the same integrating factor applied to $\rhoD$ and to
  $\rhoV$ writes each as a non-negative initial value times a positive
  exponential plus a non-negative integral of those sources. Baryons and
  radiation are unsourced and the same argument applies with no source term.
\end{proof}

Two remarks on what the theorem does and does not assert. It is conditional on
existence --- it says the densities cannot change sign while a solution lasts,
not that one lasts forever --- and existence on a maximal interval follows from
\cref{thm:regular-future} under the extra hypotheses stated there, not from
this argument. And ``forward'' is essential; see the first of the two
consequences recorded below.

The mechanism is visible in the proof and is worth naming, because it is what
the pathological models lack: \emph{every outflow is proportional to the
density of the sector it leaves}. The converse half of the dichotomy is what
makes the comparison informative, and it is equally elementary.

\begin{proposition}[Recipient-proportional couplings are not positivity-safe]
  \label{prop:recipient-proportional}
  Consider two sectors with $\dot\rho_1+3\Hub\rho_1=-Q$ and
  $\dot\rho_2+3\Hub(1+w)\rho_2=+Q$, and let $Q=\xi \Hub\rho_2$ with $\xi>0$.
  \begin{enumerate}
    \item On the face $\{\rho_1=0\}$ the field is
      $\dot\rho_1=-\xi \Hub\rho_2<0$ whenever $\rho_2>0$, so the subtangent
      condition fails on a set of codimension one and the non-negative orthant
      is \emph{not} forward invariant.
    \item Whether a given trajectory reaches that face is a separate question
      with a separate answer. Writing $u=\rho_1a^3$ gives
      $\dd u/\dd \efold=-\xi\rho_{2i}e^{-3w \efold}$, so $\rho_1$ reaches zero
      at finite $a$ for every $\xi>0$ if and only if $w\le0$, and for $w>0$ only
      if $\xi>3w\rho_{1i}/\rho_{2i}$.
  \end{enumerate}
  The same argument with the roles exchanged applies to $Q=\xi \Hub\rho_1$
  draining $\rho_2$.
\end{proposition}

The second part is not a technicality, and the first is what the comparison
rests on. A stiff recipient, $w=1$, makes the integral converge: the outflow
dies away faster than the donor dilutes, and for small enough $\xi$ the donor
survives. The failure is therefore guaranteed exactly in the regime the
catalogued models occupy --- a recipient that behaves like dark energy, $w$
close to $-1$ --- and the unqualified statement that $\rho_1$ ``is driven
through zero'' is true there and not in general.

It is worth stating the sharp criterion, because \emph{donor versus recipient}
is a useful label and not the underlying condition. By Nagumo the orthant is
forward invariant if and only if the outflow vanishes on the face it would
leave through, $Q|_{\rho_1=0}\le0$. Donor-proportionality, $Q\propto\rho_1$, is
the simplest sufficient case and the one this model realises --- and it
additionally gives uniqueness, which a coupling like $Q=\xi
\Hub\sqrt{\rho_1\rho_2}$ satisfies the invariance condition without. What the
recipient-proportional couplings fail is the condition itself:
$Q|_{\rho_1=0}=\xi \Hub\rho_2>0$, an outflow from a sector that has nothing
left. Here the precursor is a reservoir that empties, and an empty reservoir
supplies nothing.

Two consequences should be recorded immediately. First,
\cref{thm:quasi-positivity} is a statement about the future. It is not
time-symmetric: run backwards, $\rhoV$ decreases and can reach zero at a finite
scale factor. That is not a pathology of the theorem but a genuine feature of
the model, and \cref{sec:solvable-early} shows what must be assumed to avoid
it. Second, positivity here is structural rather than a consequence of tuning:
no choice of $\Gamz$, $\alpha$, $\beta$ or $\tauM$ can violate it.

\subsection{Which interacting dark energy model is this?}
\label{sec:positivity-map}

\begin{proposition}[Exact map]\label{prop:ide-map}
  With $\rhodmeff=\rhoI+\rhoD$ and $s=\rhoI/\rhodmeff$,
  %
  \begin{keyresult}{I}{eq:positivity-map}
    \QV=\xieff\,\Hub\,\rhodmeff ,
    \qquad
    \xieff=\frac{(1-\bfrac)\Gam s}{\Hub}\ \ge 0 .
  \end{keyresult}
  %
  The model reduces to the standard coupling $Q=\xi \Hub\rho_{\rm dm}$ with
  constant $\xi$ if and only if $\bfrac=0$ and $\Gam=\xi \Hub$. It can never
  take the form $Q\propto\rho_{\rm de}$.
\end{proposition}

So the model is not a relabelling of a familiar parametrization: it occupies a
restricted corner of one family, with $\xieff$ a specific decreasing function
of scale factor determined by the kinetics, and it is disjoint from the other.
This matters when comparing with published constraints, because those are
quoted for the general family.

\subsection{The prediction that cannot be fitted away}
\label{sec:positivity-no-sign-flip}

\begin{theorem}[No sign flip]\label{thm:no-sign-flip}
  Define $\delta_{\rm eff}=\dd\ln(\rhodmeff a^3)/\dd\ln a$. Under
  \cref{asm:closure} with $\rhodmeff>0$,
  %
  \begin{keyresult}{C}{eq:positivity-no-sign-flip}
    \delta_{\rm eff}(a)=-\frac{\QV}{\Hub\rhodmeff}=-\xieff(a)\le0
    \qquad\text{for every }a ,
  \end{keyresult}
  %
  for every admissible rate law, with equality only where $\Gam\rhoI=0$ or
  $\bfrac=1$. The same holds for the total dust.
\end{theorem}

\begin{proof}
  Summing the first two equations of \cref{eq:cosmology-system} gives
  $\dd(\rhodmeff a^3)/\dd t=-\QV a^3$. Divide by $\rhodmeff a^3>0$ and by
  $\Hub>0$; then $\QV=(1-\bfrac)\Gam\rhoI\ge0$ by
  \cref{thm:quasi-positivity}. Baryons are conserved and only enter the
  denominator.
\end{proof}

\begin{corollary}[Memory does not enter the sign]\label{cor:memory-invisible}
  The memory parameters $(\beta,\tauM)$ affect
  \cref{eq:positivity-no-sign-flip} only through the scalar $\Gam\ge0$. They can
  change the magnitude of $\delta_{\rm eff}$ by a bounded factor and shift the
  epoch at which it peaks, but they cannot change its sign.
\end{corollary}

We state the confrontation with data plainly, because it runs against the
model. Fits in the DESI DR2 era weakly prefer the opposite sign, that is,
energy flowing from dark energy into dark matter
\parencite{kandel2026_260712283,yang2026_260821938,wang2026_260903062}. By
\cref{thm:no-sign-flip} this is not something the one-way model can accommodate
by adjusting parameters: a robust detection of $\delta_{\rm eff}>0$ falsifies
it, independently of $\bfrac$, $\alpha$, $\beta$, $\tauM$ and the activation
profile. Two qualifications belong alongside that statement, and neither is a
defence of the model. The current preference is weak and is entangled with the
choice of coupling; and by \cref{prop:recipient-proportional} some of the
parametrizations carrying that preference are recipient-proportional and
therefore admit negative densities somewhere in their own parameter space,
which argues for repeating such analyses with a donor-proportional coupling
rather than for dismissing the result. If the
preference strengthens, the disciplined response is the reverse channel of
\cref{sec:positivity-track-b}, which is a different model, not a
reparametrization of this one.

\subsection{The expansion history is bracketed}\label{sec:positivity-sandwich}

The draft version of this work argued that the model would appear phantom to an
observer. That argument does not survive scrutiny: what a fit reports depends
on what it holds fixed. The following replaces it with a statement that is
independent of that choice.

\begin{theorem}[Sandwich]\label{thm:sandwich}
  Let the solution be defined and non-negative on $[a_e,1]$ with radiation
  exactly conserved, and define the matter deficit
  %
  \begin{keyresult}{C}{eq:positivity-deficit}
    \Delta\eqdef\frac{\rho_{\rm m}^{\rm early}-\rho_{\rm m0}}{3\Mpl^2\Hz^2}
    =\frac{1}{3\Mpl^2\Hz^2}\int_{t_e}^{t_0}\QV a^3\,\dd t\ \ge 0 ,
    \qquad
    \rho_{\rm m}^{\rm early}\eqdef\rhom(t_e)\,a_e^3 ,
  \end{keyresult}
  %
  both densities being comoving, that is, referred to $a=1$ today.
  %
  Then for all $a\in[a_e,1]$
  %
  \begin{keyresult}{C}{eq:positivity-sandwich}
    E^2_{\LCDM}\!\left(\Omega_{\mathrm{m}0},\Omega_{V0}\right)\le E^2(a)\le
    E^2_{\LCDM}\!\left(\Omega_{\mathrm{m}0}+\Delta,\ \Omega_{V0}-\Delta\right) ,
  \end{keyresult}
  %
  where both comparison models are spatially flat, carry the same radiation
  density, and share the observed $\Hz$, differing only by moving the amount
  $\Delta$ from the cosmological constant to matter. Equality holds at $a=1$ on
  both sides. If $\Delta>\Omega_{V0}$ the upper comparison model has a negative
  cosmological constant; it remains a valid bound on $E^2$, but ceases to be a
  physical $\LCDM$ cosmology.

\end{theorem}

\begin{remark}[The naive bracketing is wrong]\label{rem:sandwich-naive}
  The version of this statement one writes down first is false, so it is worth
  recording. Bracketing instead by the pair
  $(\rho_{\rm m}^{\rm early},V_{\rm early})$ does \emph{not} give an upper
  bound: that model is a lower bound, and it does not share $\Hz$.
\end{remark}

\begin{proof}
  Write $3\Mpl^2\Hub^2=m(t)a^{-3}+\rho_{\rm r0}a^{-4}+\rhoV(t)$ with
  $m=\rhom a^3$, exact by the Friedmann constraint; both comparison models
  carry the same radiation density $\rho_{\rm r0}$, which therefore cancels
  from every difference below. From \cref{eq:cosmology-system},
  $\dot m=-\QV a^3$ and $\dot \rhoV=\QV$.

  For the lower bound, integrate both backwards from today:
  $m(t)-\rho_{\rm m0}=\int_t^{t_0}\QV a'^3\,\dd t'$ and
  $\rhoV(t)-V_0=-\int_t^{t_0}\QV\,\dd t'$, so
  %
  \begin{equation}\label{eq:positivity-sandwich-lower}
    3\Mpl^2\left[\Hub^2-\Hub_{\rm low}^2\right]
    =\int_t^{t_0}\QV\left[\left(\frac{a'}{a}\right)^{3}-1\right]\dd t'\ \ge\ 0 ,
  \end{equation}
  %
  since $a'\ge a$ and $\QV\ge0$.

  The upper bound is not the same integral with the endpoints exchanged; it has
  a different structure, because the high comparison model is anchored at
  $t_e$ while $\Hz$ is fixed today. Splitting the range at $t$,
  %
  \begin{equation}\label{eq:positivity-sandwich-upper}
    3\Mpl^2\left[\Hub_{\rm high}^2-\Hub^2\right]
    =\int_{t_e}^{t}\QV a'^3\left(a^{-3}-1\right)\dd t'
    +\int_{t}^{t_0}\QV\left(1-a'^3\right)\dd t' ,
  \end{equation}
  %
  and both terms are non-negative for $a\le1$, since $a'\le1$ throughout.
  Equality holds at $a=1$ on both sides because both integrals then have
  vanishing range or vanishing integrand.
\end{proof}

The practical content is that the model cannot mimic an arbitrary expansion
history: it lives in a band whose width is one number, the matter deficit
$\Delta$, which is the difference between the matter density inferred from the
early universe and the one inferred locally. That is the honest observable, and
it is what \cref{sec:observations} reports, together with $\Hub(z)$ and
$\jerk(z)$, in place of an anchoring-dependent effective $w$.

\subsection{The reverse channel}\label{sec:positivity-track-b}

Adding a channel that returns vacuum energy to the localized sector,
%
\begin{equation}\label{eq:positivity-track-b}
  \dot \rhoV=(1-\bfrac)\Gam\rhoI-\Gamdn \rhoV ,
  \qquad
  \dot\rhoD+3\Hub\rhoD=\bfrac\Gam\rhoI+\Gamdn \rhoV ,
\end{equation}
%
preserves \cref{thm:quasi-positivity} verbatim, because the new outflow is
again proportional to the density of the sector it leaves. What it does change
is \cref{thm:no-sign-flip}: the net transfer
$\Qnet=(1-\bfrac)\Gam\rhoI-\Gamdn \rhoV$ becomes negative once the vacuum
dominates, generically at late times, so $\delta_{\rm eff}$ can flip sign. In
the four-form picture of \cref{sec:realizations-bridge} this channel is nothing
exotic: it is ordinary downward membrane nucleation, the process the upward
one competes against.

The two tracks are therefore distinguished by the sign of a single
cosmographic quantity, which is the subject of \cref{sec:cosmography}. Both
appear in the model comparison of \cref{sec:observations}.

\paragraph{What this section establishes.}
That positivity holds structurally rather than by tuning; the exact
correspondence with the standard interacting parametrizations; and a
falsifiable sign prediction that no choice of parameters can evade. What it does
not establish is that the sign predicted is the one nature chose; current data
lean the other way, weakly, and the paper reports that rather than working
around it.
